\documentclass[11pt]{article}
\usepackage[utf8]{inputenc}
\usepackage{amsmath, amssymb}
\usepackage{algorithm}
\usepackage{algpseudocode}
\usepackage{geometry}
\usepackage{booktabs}
\usepackage{tabularx}
\usepackage{xcolor}
\usepackage{hyperref}
\geometry{a4paper, margin=1in}

\definecolor{common}{RGB}{20,100,40}
\definecolor{diff}{RGB}{160,30,30}
\definecolor{gray}{RGB}{90,90,90}

\title{\textbf{PPAC Algorithm Comparison}\\[0.3em]
\large Curtin et al.\ (2010) · Existing Evaluator · Fixed Optimizer}
\date{}
\begin{document}
\maketitle

% ─────────────────────────────────────────────────────────────
\section*{Shared Notation}
% ─────────────────────────────────────────────────────────────

\begin{tabular}{@{}ll@{}}
$I,\,i$   & incident locations with weights $a_i$ \\
$J,\,j$   & candidate command-centre locations \\
$P$       & number of patrol areas to locate \\
$S$       & service distance threshold \\
$d_{ij}$  & road-network distance from incident $i$ to facility $j$ \\
$N_i$     & $\{j \in J \mid d_{ij} \le S\}$ — facilities that cover incident $i$ \\
$x_j$     & 1 if facility placed at $j$, else 0 \\
$y_i$     & 1 if incident $i$ is covered by $\ge 1$ facility \\
$O$       & maximal covering objective $= \sum_i a_i y_i$ \\
$B$       & backup objective $= \sum_i a_i \cdot (\text{count of facilities covering } i)$ \\
\end{tabular}

% ─────────────────────────────────────────────────────────────
\section*{Common PPAC Objective (all three)}
% ─────────────────────────────────────────────────────────────

{\color{common}
All three share the identical PPAC formulation from Church \& ReVelle (1974):
\[
  \text{Maximise}\;\; O = \sum_{i \in I} a_i\, y_i
\]
\[
  \text{s.t.}\quad \sum_{j \in N_i} x_j \ge y_i\;\;\forall i;\quad
  \sum_j x_j = P;\quad x_j,y_i \in \{0,1\}
\]
All three also compute the backup objective $B = \sum_i a_i \cdot y_i^{\text{count}}$,
use road-network shortest-path distances for $d_{ij}$,
and evaluate coverage via single-source Dijkstra with cutoff $S$.
}

% ─────────────────────────────────────────────────────────────
\section*{Algorithm 1 — Curtin et al.\ (2010) [PDF Paper]}
% ─────────────────────────────────────────────────────────────

\begin{algorithm}[H]
\caption{Classical PPAC (Curtin et al., 2010)}
\begin{algorithmic}[1]
\State \textbf{Input:} Geocoded incidents $I$ with Signal Code weights $a_i$;
       {\color{diff}existing beat-centroid polygons} as candidate set $J$;
       $P$ (sectors per division); $S$ (2 mi, or 1 mi for Central Division)
\State \textbf{Distance:} {\color{diff}Proprietary NCTCOG GIS network};
       all-to-all OD matrix computed externally
\State \textbf{Build} $N_i \leftarrow \{j \in J \mid d_{ij} \le S\}$ for all $i$
\State \textbf{Solve} via {\color{diff}CPLEX 8.1 branch-and-bound} (proven optimal):
\Statex \quad $\max \sum_i a_i y_i$ \; s.t.\ coverage + cardinality constraints
\State \textbf{Assign} beats to sectors from optimal $x_j^*$
\State \textbf{Evaluate} $O$, total distance; compare to hand-drawn baseline
\State \textbf{Optional backup:} relax $y_i \in \{0,\dots,P\}$;
       solve repeatedly for range of $O$ $\to$ Pareto frontier of $O$ vs $B$
\end{algorithmic}
\end{algorithm}

\noindent\textbf{Key parameters:} $|I|$ up to 87,603; $|J|$ = beat centroids within one division;
solved independently per division; $S$ empirically chosen per division.

% ─────────────────────────────────────────────────────────────
\section*{Algorithm 2 — Existing Infrastructure Evaluator}
% ─────────────────────────────────────────────────────────────

\begin{algorithm}[H]
\caption{Existing LAPD Station Coverage Evaluator}
\begin{algorithmic}[1]
\State \textbf{Input:} LAPD crime CSV ($|I|=997{,}488$, \texttt{crime\_weight} $a_i$);
       {\color{diff}21 fixed real station locations} as $J$ (no selection step);
       $S = 2.5$ mi $= 4{,}023$ m
\State \textbf{Distance:} OSMnx OSM graph projected to EPSG:4326;
       {\color{diff}no OD matrix — ego-graph Dijkstra per station at evaluation time}
\State \textbf{No facility selection:} {\color{diff}$J$ is fixed}, $x_j = 1$ for all $j$
\State \textbf{For each station $j \in J$:}
\Statex \quad Run Dijkstra$(G, j, \text{cutoff}=S)$ $\to$ reachable node set
\Statex \quad Mark all incidents whose OSM node is reachable: $y_i \mathrel{+}= 1$
\State \textbf{Compute} $O = \sum_i a_i \cdot \mathbf{1}[y_i \ge 1]$;
       $B = \sum_i a_i \cdot y_i$
\State \textbf{Output:} coverage statistics + map; {\color{diff}no optimisation performed}
\end{algorithmic}
\end{algorithm}

\noindent\textbf{Key parameters:} $|J|=21$ (fixed); city-wide single pass;
$S=2.5$ mi; $a_i =$ \texttt{crime\_weight}; no candidate generation step.

% ─────────────────────────────────────────────────────────────
\section*{Algorithm 3 — Fixed Greedy Optimizer}
% ─────────────────────────────────────────────────────────────

\begin{algorithm}[H]
\caption{Fixed PPAC Optimizer (Greedy + Road-Network)}
\begin{algorithmic}[1]
\State \textbf{Input:} LAPD crime CSV ($|I|=997{,}488$, \texttt{crime\_weight} $a_i$);
       city boundary GeoJSON; $P=21$; $S=2.5$ mi
\State \textbf{Stage 1 — Candidate generation} {\color{diff}(data-driven, replaces existing beat polygons):}
\Statex \quad Weighted Mini-Batch K-Means on incidents $\to$ 300 beat centroids $= J$
\State \textbf{Distance:} {\color{diff}Two OSMnx graph projections:}
\Statex \quad $G_\text{4326}$ for vectorised node snapping (KD-tree in degree space)
\Statex \quad $G_\text{metric}$ (UTM) for Dijkstra (edge lengths unambiguously in metres)
\State \textbf{Build coverage sets} $C_j \leftarrow \{i \in I \mid d_{ij} \le S\}$ for all $j \in J$
       via Dijkstra$(G_\text{metric}, j, \text{cutoff}=S)$
\State \textbf{Stage 2 — Greedy PPAC} {\color{diff}(replaces K-Means clustering of beat centres):}
\Statex \quad $\text{covered} \leftarrow \emptyset$; \; $H \leftarrow \emptyset$
\For{$k = 1$ \textbf{to} $P$}
  \State $j^* \leftarrow \arg\max_{j \notin H}\; \sum_{i \in C_j \setminus \text{covered}} a_i$
  \State $H \leftarrow H \cup \{j^*\}$; \quad $\text{covered} \leftarrow \text{covered} \cup C_{j^*}$
\EndFor
\State \textbf{Compute} $O = \sum_i a_i \cdot \mathbf{1}[i \in \text{covered}]$;
       $B = \sum_i a_i \cdot |\{j \in H : i \in C_j\}|$
\State \textbf{Output:} Voronoi beat map + sector map; coverage CSV
\end{algorithmic}
\end{algorithm}

\noindent\textbf{Key parameters:} $|J|=300$ (derived); city-wide single pass;
$P=21$; $S=2.5$ mi; greedy gives $(1-1/e)\approx 63\%$ approximation guarantee on $O$.

% ─────────────────────────────────────────────────────────────
\section*{Differences at a Glance}
% ─────────────────────────────────────────────────────────────

\renewcommand{\arraystretch}{1.35}
\begin{tabularx}{\textwidth}{@{} l X X X @{}}
\toprule
\textbf{Aspect} & \textbf{PDF (Curtin)} & \textbf{Evaluator} & \textbf{Fixed Optimizer} \\
\midrule
\textbf{Purpose} & Optimize geography & Audit fixed stations & Optimize geography \\
\textbf{Candidate set $J$} & Existing beat centroids & 21 real stations (fixed) & 300 K-Means centroids \\
\textbf{Facility selection} & IP (proven optimal) & None — all selected & Greedy ($\ge 63\%$ opt.) \\
\textbf{$P$} & Per-division (e.g.\ 5) & 21 (all used) & 21 \\
\textbf{$S$} & 2 mi (1 mi Central) & 2.5 mi & 2.5 mi \\
\textbf{$a_i$} & Signal Code 1--5 & \texttt{crime\_weight} & \texttt{crime\_weight} \\
\textbf{Road network} & Proprietary NCTCOG & OSMnx EPSG:4326 & OSMnx UTM + 4326 \\
\textbf{OD matrix} & All-to-all (external GIS) & None (ego-graph only) & Per-candidate ego-graph \\
\textbf{Backup coverage} & Formal multi-obj.\ IP & Computed post-hoc & Computed post-hoc \\
\textbf{Scope} & Per division & City-wide & City-wide \\
\textbf{Solver} & CPLEX 8.1 & N/A & Greedy heuristic \\
\bottomrule
\end{tabularx}

\bigskip
\noindent\textbf{Why the evaluator outscored the original optimizer.}
The original optimizer ran K-Means Stage 2 on beat centroids, minimizing Euclidean
variance — an objective unrelated to $O$. The 21 real LAPD stations were
historically sited near high-crime areas, so they outperformed geometrically-centered
K-Means HQs on the coverage metric $O$. The fixed optimizer replaces Stage 2
with the greedy algorithm that directly maximizes $O$, and uses $G_\text{metric}$
for Dijkstra so that the road-distance cutoff is exact.

\end{document}